\section{Comparison of numerical and fixed-ray psh functions on an ample ray}
\label{sec:native-numerical-psh-ray-comparison}

Let $p:X\to\operatorname{Spec}k$ be projective with $X$ integral, let $L$ be
an ample invertible sheaf and $a\ge0$.  A fixed-ray psh function of $(L,a)$ is
an upper semicontinuous function, not identically $-\infty$, which is the
pointwise limit of a decreasing net of functions in $\mathcal H(L,a_j)$ with
$a_j\in\mathbb Q$ and $a_j\to a$.  By BJ22, Definition~4.1, a function of
$\operatorname{PSH}(X,a[L])$ is instead a decreasing limit of varying-line
Fubini--Study functions whose classes converge to $a[L]$ in $N^1(X)$.  This
section proves that the two classes coincide.  The argument uses two standard
results, recorded as hypotheses of the general statements; the second is
proved below from the nef-plus-ample theorem:
\begin{itemize}
\item[(FD)] the real numerical space $N^1(X)$ is finite dimensional;
\item[(K)] Kleiman's criterion: an invertible sheaf whose numerical class
lies in the ample cone is ample.
\end{itemize}

\begin{lemma}[The ample cone is open]
\label{lem:native-numerical-ample-cone-open-fd}
\lean{EnvelopeProperty.NativeRealNumericalClasses.isOpen_ampleCone_of_finiteDimensional}
\leanok
\uses{thm:native-numerical-line-family-maps,thm:native-numerical-curve-weak-vector-topology,thm:native-numerical-ample-family-openness}
Assume (FD).  Then the ample cone is open in $N^1(X)$.
\end{lemma}
\begin{proof}
Choose a basis $(e_i)$ of $N^1(X)$ and write each $e_i$ as a finite real
combination of classes of invertible sheaves.  Collect all these sheaves into
one finite family $(L_s)$, with coefficient map $F:\mathbb R^S\to N^1(X)$.
Each $[L_s]$ lies in the range of $F$, hence so does each $e_i$, so $F$ is
surjective.  The preimage under $F$ of the ample cone is open
(the ample cone has a positive perturbation radius along every finite line
family), and a surjective linear map onto a finite dimensional Hausdorff
topological vector space is open.  Hence the ample cone, the image of its
preimage, is open.
\end{proof}

\begin{definition}[Kleiman's criterion]
\label{def:native-numerical-kleiman-criterion}
\lean{EnvelopeProperty.NativeNumericalFubiniStudy.NumericalAmplenessCriterion}
\leanok
\uses{def:native-numerical-ample-cone}
We say that (K) holds for $p$ if every invertible sheaf on $X$ whose
numerical class lies in the ample cone is ample.
\end{definition}

\begin{theorem}[Kleiman's criterion from nef plus ample]
\label{thm:native-numerical-kleiman-from-nef-ample}
\lean{EnvelopeProperty.NativeNumericalFubiniStudy.numericalAmplenessCriterion}
\leanok
\uses{def:native-numerical-kleiman-criterion,thm:native-nef-plus-ample,thm:native-numerical-positive-ample-sums,thm:native-numerical-rational-ample-lines,thm:native-numerical-ample-classes,thm:fc-original-ample-line-positive-powers,thm:native-numerical-generated-line-sums,lem:native-numerical-line-tensor-power}
If $X$ is integral, then (K) holds.
\end{theorem}
\begin{proof}
Let $T$ be invertible with $[T]$ in the ample cone.  Write
$[T]=\sum_ic_i[A_i]$ with finitely many ample $A_i$ and $c_i>0$, choose
rationals $q_i$ with $c_i/2<q_i<c_i$, and choose an integer $d>0$ and an ample
$B$ with $\sum_iq_i[A_i]=d^{-1}[B]$.  The sheaf $N=T^{\otimes d}\otimes B^\vee$
has class $d[T]-[B]=\sum_id(c_i-q_i)[A_i]$, a positive combination of ample
classes, hence nef.  By the nef-plus-ample theorem $N\otimes B$ is ample; it is
isomorphic to $T^{\otimes d}$ by evaluation $B^\vee\otimes B\cong\mathcal O_X$.
A line with an ample positive power is ample, so $T$ is ample.
\end{proof}

\begin{lemma}[Membership in a fixed ray from an ample difference class]
\label{lem:native-numerical-ray-from-ample-class}
\lean{EnvelopeProperty.NativeNumericalFubiniStudy.FSData.isFS_ray_of_mem_ampleCone}
\leanok
\uses{def:native-numerical-kleiman-criterion,cor:native-numerical-ray-realization-difference}
\source{bj22-trivval:prop:FS}
Assume (K).  Let $D=(M,m,\dots)$ be a nonzero varying-line datum and let
$b\ge0$ be rational with $b[L]-[D]$ in the ample cone.  Then
$\varphi_D\in\mathcal H(L,b)$.
\end{lemma}
\begin{proof}
Write $b=u/v$ with $u\ge0$ and $v>0$ integers, and put $r=v$, $d=mu$.  The
sheaf $A=L^{\otimes d}\otimes(M^{\otimes r})^\vee$ has class
$mr(\frac{d}{mr}[L]-[D])=mv(b[L]-[D])$, a positive multiple of an ample class,
so $A$ is ample by (K).  By
Corollary~\ref{cor:native-numerical-ray-realization-difference},
$\varphi_D\in\mathcal H(L,d/(mr))=\mathcal H(L,b)$.
\end{proof}

\begin{lemma}[A ray coefficient for every datum]
\label{lem:native-numerical-ray-coefficient}
\lean{EnvelopeProperty.NativeNumericalFubiniStudy.FSData.exists_ray_coefficient}
\leanok
\uses{def:native-numerical-whole-section-fs-data,thm:native-numerical-ample-family-openness,lem:native-numerical-ample-cone-nef}
Let $L$ be ample and $D=(M,m,\dots)$ a varying-line datum.  Then
$t[L]-[D]$ is ample for some $t>0$.
\end{lemma}
\begin{proof}
The ample class $[L]$ has a positive perturbation radius $\varepsilon$ along
the one-line family $(M)$, so $[L]-\frac\varepsilon2[M]$ is ample.  With
$t=(\frac\varepsilon2m)^{-1}$ we get
$t[L]-[D]=t([L]-\frac\varepsilon2[M])$, a positive multiple of an ample class.
\end{proof}

\begin{theorem}[The net argument on an ample ray]
\label{thm:native-numerical-psh-ray-net}
\lean{EnvelopeProperty.NativeNumericalFubiniStudy.isPSH_ray_of_antitone_net}
\leanok
\uses{lem:native-numerical-ray-coefficient,lem:native-numerical-ray-from-ample-class,def:native-line-whole-section-psh}
\source{bj22-trivval:thm:pshample}
Assume (K), let $L$ be ample and $a\ge0$.  Let $\varphi$ be upper
semicontinuous, not identically $-\infty$, and the pointwise infimum of the
decreasing functions of nonzero varying-line data $D_j$, $j\in J$ (a nonempty
directed preorder), such that for every $\varepsilon>0$ the class
$(a+\varepsilon)[L]-[D_j]$ is ample for all $j$ large enough.  Then $\varphi$
is a fixed-ray psh function of $(L,a)$.
\end{theorem}
\begin{proof}
Let $\ell=[L]$.  For each $j$ let $T_j$ be the set of real $t\ge a$ with
$t\ell-[D_j]$ ample.  If $s\in T_j$ and $t>s$, then
$t\ell-[D_j]=(s\ell-[D_j])+(t-s)\ell$ is ample, so $T_j$ is upward closed.  By
Lemma~\ref{lem:native-numerical-ray-coefficient} and upward closedness, $T_j$
is nonempty; it is bounded below by $a$.  Let $\sigma_j=\inf T_j$.  For
$\varepsilon>0$ the hypothesis gives $a+\varepsilon\in T_j$, hence
$\sigma_j\le a+\varepsilon$, for all $j$ large enough.

Index a new net by $J\times\mathbf N$ with the product order, which is
directed.  For $(j,n)$ choose a rational $b_{j,n}$ with
$\sigma_j<b_{j,n}<\sigma_j+\frac1{n+1}$.  Some $t\in T_j$ satisfies
$t<b_{j,n}$, so $b_{j,n}\in T_j$; in particular $b_{j,n}\ge a\ge0$ and
$b_{j,n}\ell-[D_j]$ is ample, so $\varphi_{D_j}\in\mathcal H(L,b_{j,n})$ by
Lemma~\ref{lem:native-numerical-ray-from-ample-class}.  The functions
$\varphi_{D_j}$, indexed by $(j,n)$, decrease and have infimum $\varphi$.
Finally $b_{j,n}\to a$: always $b_{j,n}\ge a$, and given $c>a$ put
$\varepsilon=(c-a)/2$, choose $j_0$ with $\sigma_j\le a+\varepsilon$ for
$j\ge j_0$ and $n_0$ with $\frac1{n_0+1}<\varepsilon$; for
$(j,n)\ge(j_0,n_0)$ we get $b_{j,n}<a+\varepsilon+\frac1{n+1}<c$.
\end{proof}

\begin{theorem}[Comparison on an ample ray]
\label{thm:native-numerical-psh-ray-comparison}
\lean{EnvelopeProperty.NativeNumericalFubiniStudy.isPSH_ray_of_isCurveWeakPSH,EnvelopeProperty.NativeNumericalFubiniStudy.isCurveWeakPSH_iff_isPSH_ray}
\leanok
\uses{lem:native-numerical-ample-cone-open-fd,thm:native-numerical-psh-ray-net,def:native-numerical-curve-weak-psh,thm:native-numerical-curve-weak-psh-from-ray}
\source{bj22-trivval:defi:psh,bj22-trivval:thm:pshample}
Assume (FD) and (K), let $L$ be ample and $a\ge0$.  A function is curve-weak
numerical psh of class $a[L]$ if and only if it is a fixed-ray psh function
of $(L,a)$.
\end{theorem}
\begin{proof}
A fixed-ray psh function is curve-weak numerical psh of class $a[L]$
(Theorem~\ref{thm:native-numerical-curve-weak-psh-from-ray}).  Conversely, let
$\varphi$ be the decreasing limit of the functions of nonzero data $D_j$
with $[D_j]\to a[L]$.  For $\varepsilon>0$ the classes
$(a+\varepsilon)[L]-[D_j]$ converge to the ample class $\varepsilon[L]$, and
the ample cone is open (Lemma~\ref{lem:native-numerical-ample-cone-open-fd}),
so they are eventually ample.  Apply
Theorem~\ref{thm:native-numerical-psh-ray-net}.
\end{proof}

\begin{corollary}[Comparison on an ample ray, assuming only (FD)]
\label{cor:native-numerical-psh-ray-comparison-fd}
\lean{EnvelopeProperty.NativeNumericalFubiniStudy.isCurveWeakPSH_iff_isPSH_ray'}
\leanok
\uses{thm:native-numerical-psh-ray-comparison,thm:native-numerical-kleiman-from-nef-ample}
Assume (FD), let $L$ be ample and $a\ge0$.  A function is curve-weak numerical
psh of class $a[L]$ if and only if it is a fixed-ray psh function of $(L,a)$.
\end{corollary}
\begin{proof}
Theorem~\ref{thm:native-numerical-psh-ray-comparison}, with (K) supplied by
Theorem~\ref{thm:native-numerical-kleiman-from-nef-ample}.
\end{proof}

\begin{remark}
Combined with the comparison of fixed-ray and model psh functions for the
model given by an embedding with $\mathcal O(1)|_X\cong L^{\otimes r}$, this
identifies BJ22's $\operatorname{PSH}(X,a[L])$ with the model psh class of
coefficient $a/r$, under (FD).
\end{remark}
